\section{Related Work}\label{sec:related}

\scl{} touches four research directions: the trackers it serves, the
confidence of the detectors that feed them, trackers that adapt their own
thresholds, and methods that adapt inference at run time. They are reviewed
below in this order, and Table~\ref{tab:related} compares the closest
methods with \scl{}.

Most practical trackers follow the tracking-by-detection design. SORT
predicts each box with a Kalman filter and matches the predictions to new
detections by overlap with the Hungarian method \cite{bewley2016sort}.
ByteTrack introduced a second association stage that uses low-confidence
detections to continue existing tracks \cite{zhang2022bytetrack}, and many
later trackers kept this two-stage design. Others, such as OC-SORT
\cite{cao2023ocsort}, use a single confidence threshold. All of these
trackers read detector scores through fixed thresholds that were tuned for
one detector, usually YOLOX \cite{ge2021yolox}. Seven published trackers of
both kinds are used in this paper as hosts and are kept unchanged
(Sect.~\ref{sec:setup}).

These fixed thresholds are a weak point, because detector scores do not
carry the same meaning across models. Detectors of different families, such
as YOLOX \cite{ge2021yolox}, YOLOv8 \cite{jocher2023yolov8}, YOLOv9
\cite{wang2024yolov9}, DETR \cite{carion2020detr}, and RT-DETR
\cite{zhao2024rtdetr}, produce different score distributions. Their
confidence is not calibrated across models and conditions, and calibration
methods for object detection fit a mapping on labelled data
\cite{kuppers2020detcalib}. A tracker tuned for one detector therefore
cannot assume that another detector's scores mean the same thing, and labels
to recalibrate them are rarely available after deployment. A label-free tool
for this exists in image processing: Otsu's method selects a threshold from
a histogram by maximizing the separation between two classes
\cite{otsu1979}. \scl{} applies it to the logits (log-odds) of the detector
scores of recent frames.

Some trackers already adapt their own confidence thresholds. Adaptive
ByteTrack derives the confidence threshold of ByteTrack from the detections
of each frame \cite{ma2023adaptivebytetrack}. AdapTrack adapts the
thresholds of its matching step \cite{shim2024adaptrack}. ACT sets the high
and low thresholds of each frame from the first quartile of the confidence
scores and the number of candidates \cite{act2026adaptive}. These methods
are built into one tracker and use statistics of the current frame. \scl{}
differs in four ways. Its statistics come from past frames. It decides for
the whole stream whether to intervene at all, and otherwise keeps the
tracker's own operating point. It acts through a declared host contract on
several published trackers without modifying them. It is evaluated as one
frozen policy under confidence shift and on external systems.

A wider family of methods adapts inference itself. Dynamic neural networks
adapt their computation to each input \cite{han2022dynamic}, and Chameleon
re-profiles the configuration of video analytics pipelines over time
\cite{jiang2018chameleon}. These methods learn or profile a policy for each
model. Our earlier work, AC-MOT \cite{ismail2026acmot}, is a training-free
controller that sets the detector operating point (score threshold, NMS
threshold, and input resolution) online from image and track cues. Such a
controller is tied to the detector it was calibrated with, because a change
of the detector's score scale leaves the image unchanged. \scl{} reads the
score stream instead and acts on the tracker's input and thresholds.

In summary, existing methods either adapt thresholds inside one tracker from
the current frame, calibrate scores with labelled data, or learn a policy for
one model. A training-free layer that calibrates itself from the detector's
own score stream, serves many unmodified trackers, and keeps a well-matched
tracker at its own accuracy is missing. \scl{} is designed to fill this gap.
Table~\ref{tab:related} summarizes the differences.

\begin{table}[t]
\caption{Summary of the adaptive methods closest to \scl{} (BT: ByteTrack).}
\label{tab:related}
\centering\footnotesize
\setlength{\tabcolsep}{3pt}
\begin{tabular}{@{}>{\raggedright\arraybackslash}p{17mm}>{\raggedright\arraybackslash}p{18mm}>{\raggedright\arraybackslash}p{15mm}>{\raggedright\arraybackslash}p{18mm}@{}}
\toprule
Method & Adapted & Statistics & Where\\
\midrule
Adapt.\ BT \cite{ma2023adaptivebytetrack} & confidence threshold & current frame & inside BT\\
AdapTrack \cite{shim2024adaptrack} & matching thresholds & current frame & inside the tracker\\
ACT \cite{act2026adaptive} & high and low thresholds & current frame & inside the tracker\\
K{\"u}ppers \cite{kuppers2020detcalib} & detector scores & labelled data & detector output\\
Chameleon \cite{jiang2018chameleon} & pipeline configuration & profiling & analytics pipeline\\
AC-MOT \cite{ismail2026acmot} & detector operating point & image and track cues & before the detector\\
\scl{} (this paper) & candidates, scores, thresholds & past frames, no labels & between detector and tracker\\
\bottomrule
\end{tabular}
\end{table}
